\documentclass[12pt]{article}

\usepackage{geometry}
\usepackage{hyperref}
\usepackage{fancyhdr}
\usepackage{enumitem}
\usepackage{setspace}
\usepackage{parskip}


\geometry{a4paper, margin=1in}

\pagestyle{fancy}
\fancyhf{}
\fancyhead[L]{Linux} 
\fancyfoot[C]{\thepage}

\title{Basics Linux Commands}
\author{0xSaad / Saad Almalki}
\date{\today}

\begin{document}

\maketitle

\section{Introduction to commands}

There are many simple and basic commands such as :
\begin{enumerate}
    \item \textbf{ls} Using to show a list of all files and directories in same directory.
    \item \textbf{cd} Using to enter another directory.
\end{enumerate}

And this commands can be used to other jobs with more optionals :

\begin{enumerate}
    \item \textbf{ls -l} Showing a long list of all files and directories and more details about it.
    \item \textbf{ls -a} Showing a list of all files/directories with hidden and secret files and directories.
    \item \textbf{ls -la} "ls -l" + "ls -a"
    \item \textbf{cd ..} Back to the parent directory (the parent of current directory/folder)
    \item \textbf{cd -} Back to the previous directory , You were in before the current on.
\end{enumerate}



\section{Create/Delete commands}
\begin{enumerate}
    \item \textbf{mkdir} Make a new directory
    \item \textbf{rmdir} Remove a directory (It's required to be an empty directory).
    \item \textbf{rm -r} Remove a directory and all files/directories on this directory.
    \item \textbf{touch} Make a new file with file extension "touch main.py" -> Make a new python source code file
    \item \textbf{nano} Editing an exist file
    \item \textbf{rm} Remove a file
\end{enumerate}



\begin{flushright}
Created by: Saad Almalki \\
\end{flushright}

\end{document}
